\documentclass[12pt]{article}
\usepackage[utf8]{inputenc}
\usepackage[T1]{fontenc}
\usepackage{hyperref}
\title{Samberigi Language Resources for K-12 Curriculum Development}
\author{Compiled by Claude Sonnet 5.5 (claude-sonnet-5-5) for Andr\'as Kornai}
\date{}
\begin{document}
\maketitle

\section{Introduction}
Samberigi (also called Sau, ISO 639-3: ssx) is spoken in the Samberigi Valley in Papua New Guinea's Southern Highlands Province, near the Gulf border. Wikipedia gives 3,130 speakers, based on a 1981 count. The documentation is sparse. The name Sau and some other alternate names overlap with other languages, so check which language a source means.

\section{Primary Educational Resources}
No primers, textbooks or school programmes for Samberigi turned up in the search.

\section{Linguistic and Reference Materials}
\begin{itemize}
\item Wikipedia, ``Samberigi language'': an overview. \url{https://en.wikipedia.org/wiki/Samberigi_language}
\item New Guinea World profile for Samberigi (Sau). It summarises a 1973 phonology analysis by MacDonald. \url{https://newguineaworld.linguistik.uzh.ch/families/central-new-guinea-highlands/enga-southern-highlands/huli-southern-highlands/southern-highlands/sau}
\item Max Planck counting-systems page for Samberigi, with numeral data. \url{https://lingweb.eva.mpg.de/channumerals/Samberigi.htm}
\end{itemize}
The search found no full grammar or dictionary.

\section{Religious and Community Resources}
None were found apart from the audio listed below.

\section{Digital and Audio Resources}
\begin{itemize}
\item Global Recordings Network: audio Bible stories and evangelistic messages in Samberigi. \url{https://globalrecordings.net/en/language/ssx}
\end{itemize}

\section{Curriculum Development Context}
Little is documented, so a Samberigi curriculum would need new community-based work. That work would cover an orthography, word lists and local readers. The Global Recordings audio may be a starting point for spoken-language material. Any plan should be checked with Samberigi speakers and with PNG education authorities.

\end{document}